\chapter{Symbolübersicht}
\label{app:symbol-reference}

Dieser Anhang erklärt die im Regelwerk verwendeten Symbole. Ihre Bedeutung bleibt unabhängig von ihrer konkreten grafischen Gestaltung gleich.

\section{Würfel}

\begin{hptable}{C{2.2cm} Y}
\textbf{Symbol} & \textbf{Bedeutung} \\
\heroDie & Heldenwürfel. Jede Probe eines Helden enthält genau einen Heldenwürfel. \\
\greenDie & Grüner Attributswürfel. Er steht für die natürlichen Eigenschaften eines Helden. \\
\blueDie & Blauer Fertigkeitswürfel. Er steht für Übung, Wissen und Erfahrung in einer Fertigkeit. \\
\orangeDmgDie & Oranger Schadenswürfel für leichten Schaden. \\
\redDmgDie & Roter Schadenswürfel für mittleren Schaden. \\
\blackDmgDie & Schwarzer Schadenswürfel für schweren Schaden. \\
\end{hptable}

\section{Würfelergebnisse und Ressourcen}

\begin{hptable}{C{2.2cm} Y}
\textbf{Symbol} & \textbf{Bedeutung} \\
\star & Stern. Sterne zeigen Erfolge bei Proben und Angriffen. \\
\moon & Mond. Jeder gewürfelte Mond erzeugt sofort einen Mondkristall. \\
\crystal & Mondkristall. Gib ihn aus, um alle Würfel einer Farbe neu zu würfeln oder einen anderen Effekt auszulösen. Ein Held kann beliebig viele Mondkristalle besitzen. Wann ungenutzte Mondkristalle verfallen, bestimmt die Spielleitung. \\
\end{hptable}

\section{Leben und Mana}

\begin{hptable}{C{2.2cm} Y}
\textbf{Symbol} & \textbf{Bedeutung} \\
\lifeLoss & Lebenspunkte verlieren. Die angegebene Zahl wird von den aktuellen Lebenspunkten abgezogen; Lebenspunkte sinken nie unter 0. \\
\lifeGain & Lebenspunkte erhalten. Die angegebene Zahl wird zu den aktuellen Lebenspunkten addiert, höchstens bis zum Maximum. \\
\manaLoss & Mana bezahlen oder verlieren. Die angegebene Zahl wird vom aktuellen Mana abgezogen. \\
\manaGain & Mana erhalten. Die angegebene Zahl wird zum aktuellen Mana addiert, höchstens bis zum Manavorrat. \\
\end{hptable}

\section{Aktionen}

\begin{hptable}{C{2.2cm} Y}
\textbf{Symbol} & \textbf{Bedeutung} \\
\standardAction & Standardaktion. Sie wird für Angriffe, Zauber und andere größere Handlungen verwendet. \\
\moveAction & Bewegungsaktion. Sie wird für Bewegung oder das Neuordnen von Ausrüstung verwendet. \\
\freeAction & Freie Aktion. Sie steht für kurze Handlungen, die kaum Zeit kosten. \\
\end{hptable}

\section{Angriffsfertigkeiten}

\begin{hptable}{C{2.2cm} Y}
\textbf{Symbol} & \textbf{Bedeutung} \\
\meleeAttack & Schlagen. Kennzeichnet Angriffe, die mit der Fertigkeit \textbf{Schlagen} ausgeführt werden. \\
\rangedAttack & Schießen. Kennzeichnet Angriffe, die mit der Fertigkeit \textbf{Schießen} ausgeführt werden. \\
\magicAttack & Zaubern. Kennzeichnet Angriffe, die mit der Fertigkeit \textbf{Zaubern} ausgeführt werden. \\
\end{hptable}

\section{Kampf und Karten}

\begin{hptable}{C{2.2cm} Y}
\textbf{Symbol} & \textbf{Bedeutung} \\
\armor & Schild-Symbol. Jedes Schild-Symbol verhindert 1 Schaden. Schild-Symbole können durch Rüstungen, Schilde, Talente und andere Effekte gewährt werden. \\
\armorPierce & Rüstungsdurchdringung. Die angegebene Anzahl an Schild-Symbolen des Ziels wird ignoriert. \\
\evade & Ausweichen. Das Symbol kennzeichnet Ausweichproben oder Effekte, mit denen Schadenswürfel entfernt werden können. \\
\exhaust & Erschöpfen. Drehe die Karte um. Solange sie erschöpft ist, können ihre Angaben und Effekte nicht genutzt werden. \\
\end{hptable}